%! TeX root = ./main.tex
% vim: set spelllang=en:
\section{padapto --- Tools for algebraic dynamic programming}
\programdetails{https://github.com/matteodelabre/padapto}{GPL v3.0}{Between May~2025 and April~2026, 1\,100 lines of code}

The algebraic approach to dynamic programming, presented in \cref{chap:adp}, makes it possible to develop combinatorial optimization algorithms in a modular way.
\emph{padapto} is a library written in the Python language that offers the tools needed to adopt this technique for programming algorithms.

This technique is based on the formal separation between the specification of the inductive structure of the solution space considered by the algorithm and the definition of the optimization criteria through an algebra.
These two aspects are linked by the notion of signature, which defines the set of operators used on the one hand to construct the solutions and on the other hand to assign them a value.

The following modules make it possible to define signatures and evaluation algebras.
Their use is illustrated in \cref{lst:logiciel:padapto-sig-alg}.
\begin{itemize}
    \item \textbf{\texttt{padapto.signature}} defines \texttt{Signature}, a base class for declaring the signature of a problem (see \cref{sec:adp:eval}).
        This class predefines two operators common to all signatures, the nullary operator \texttt{null} and the binary operator \texttt{choose}, respectively representing the absence of a solution and the choice between two sets of solutions.
    \item \textbf{\texttt{padapto.evaluation}} contains functions for defining algebras, notably \texttt{count}, \texttt{trace}, \texttt{additive} and \texttt{boltzmann}, as well as functions for combining and transforming them (see \cref{sec:adp:alg-combin}), notably \texttt{join}, \texttt{lex}, \texttt{power}, \texttt{group}, \texttt{limit} and \texttt{pareto}.
\end{itemize}

\begin{listing}[htb]
\small
\begin{minted}{python}
@dataclass(frozen=True)
class AlignSignature[T](Signature[T]):
  unit: Callable[[], T]
  match: Callable[[T, str, str], T]
  delete: Callable[[T, str], T]
  insert: Callable[[T, str], T]

counter = count(AlignSignature)
min_unit = additive(AlignSignature, choose="min",
  match=lambda sym1, sym2: 1 if sym1 != sym2 else 0,
  delete=lambda _: 1, insert=lambda _: 1)
count_min = join(cost=min_unit, count=counter) | lex("cost")
\end{minted}
    \caption[Definitions of a signature and of algebras with the \emph{padapto} library]{
        Definitions of a signature and of algebras for the sequence alignment problem introduced in \cref{ex:seq-align}.
        The signature adds to the predefined operators \texttt{null} and \texttt{choose} the nullary operator \texttt{unit}, representing an empty alignment, and the unary operators \texttt{match}, \texttt{delete} and \texttt{insert}, representing the three possible operations in an alignment.
        The \texttt{counter} algebra counts the solutions of the solution space,
        the \texttt{min\_unit} algebra computes the best cost of an alignment by assigning a unit cost to each operation,
        while \texttt{count\_min} combines the two algebras by counting the solutions that achieve the minimum cost.
        \label{lst:logiciel:padapto-sig-alg}
    }
\end{listing}

The \texttt{trace} function of the \texttt{padapto.evaluation} module makes it possible to automatically derive the trace algebra (see \cref{sec:adp:solution-spaces:trace}) associated with a given signature.
This algebra automatically constructs the algebraic circuits representing the solution spaces associated with each instance of a problem.

Once an algebraic circuit has been obtained, it becomes possible to perform various operations on it without having to run the algorithm again.
For example, it is possible to extract solutions from it, by enumeration or by uniform or Boltzmann sampling, to evaluate other algebras on it, or to compute the frequencies of occurrence of certain nodes (see \cref{sec:adp:solution-spaces}).
These functions are provided by the following module.

\begin{itemize}
    \item \textbf{\texttt{padapto.circuit}} offers functions for manipulating algebraic circuits. \texttt{serialize} and \texttt{unserialize} make it possible to convert them to a textual format, \texttt{render} to display them, \texttt{enumerate\_solutions} and \texttt{get\_solution} to enumerate their solutions, \texttt{eval}, \texttt{eval\_inside} and \texttt{eval\_outside} to evaluate an algebra, and \texttt{sample} to sample solutions from them according to a distribution defined by an algebra.
\end{itemize}

The representation and manipulation of algebraic circuits, which are directed acyclic graphs, are done through the \emph{sowing} library presented previously.

The previous modules can be used in algorithms defined in a \enquote{traditional} way, by iterative or recursive functions specialized for a given problem.
An alternative way to solve these problems is through the yield grammars presented in \cref{sec:adp:yield-gram}.

\begin{itemize}
    \item \textbf{\texttt{padapto.structure.grammar}} allows grammar clauses to be defined in Python and implements the recursive parsing \cref{alg:seq-yield-parser} for grammars.
    \item \textbf{\texttt{padapto.structure.pattern}} defines a domain-specific sublanguage for expressing the patterns of sequences, sets, trees, forests and numbers that are used in the heads of clauses and make it possible to decompose the inputs of a dynamic programming algorithm.
\end{itemize}

\begin{listing}[p]
\small
\begin{minted}{python}
@grammar
class AlignGrammar[T](Grammar[T]):
  alg: AlignSignature[T]

  @predicate
  def align(left: str, right: str) -> T:
    return  # type: ignore

  @clause(left=Empty(), right=Empty())
  def _empty(self):
    return self.alg.unit()

  @clause(
    left=chain(Item(Var("a")), Subseq(Var("X"))),
    right=chain(Item(Var("b")), Subseq(Var("Y"))),
  )
  def _match(self, a: str, b: str, X: str, Y: str) -> T:
    return self.alg.match(self.align(left=X, right=Y), a, b)

  @clause(
    left=chain(Item(Var("a")), Subseq(Var("X"))),
    right=Var("Y"),
  )
  def _delete(self, a: str, X: str, Y: str) -> T:
    return self.alg.delete(self.align(left=X, right=Y), a)

  @clause(
    left=Var("X"),
    right=chain(Item(Var("b")), Subseq(Var("Y"))),
  )
  def _insert(self, b: str, X: str, Y: str) -> T:
    return self.alg.insert(self.align(left=X, right=Y), b)
\end{minted}
    \caption[Definition of a grammar with the \emph{padapto} library]{
        Implementation of the grammar of \cref{ex:seq-align} in Python using \emph{padapto}.
        A grammar is a class descending from \texttt{Grammar} and decorated with \texttt{@grammar}.
        The class contains a method decorated with \texttt{@predicate} for each \emph{predicate} of the grammar, each argument of which corresponds to a dimension of the predicate.
        It also contains a method  for each \emph{clause} of the grammar.
        Clauses are decorated with \texttt{@clause}, giving a decomposition pattern for each of the dimensions of the associated predicate.
        \label{lst:logiciel:padapto-gram}
    }
\end{listing}

Each class representing a grammar, such as the one defined in \cref{lst:logiciel:padapto-gram}, can be instantiated by giving it an algebra to use.
Each instance thus created holds the memoization tables that keep in memory the partial results of the subproblems.

The methods associated with each predicate are replaced by an automatically generated implementation that follows the principle of \cref{alg:seq-yield-parser}.
When a predicate is called, this method iterates through the clauses corresponding to the predicate and through the possible decompositions of the inputs on each dimension, then collects the results to store them in the associated entry of the memoization table.

Regardless of how the structure of the solution space is defined, there are two options for evaluating an algebra on it.
The first is to pass the algebra directly to the algorithm or to the grammar, which has the effect of producing a single value corresponding to the result of the evaluation.
Generally, the algebra only considers a limited part of the solution space; thus, once the result is obtained, there is no longer enough information to go through other associated solutions.
The other option is to use a trace algebra to extract the algebraic circuit describing the solution space.
